% From Meal Prep.tex. Appliance: Instant Pot (5.7 L)
\begin{recipebox}{Beef \& Mixed Bean Chilli}{6}{Prep 10 min, cook 15 min pressure + 10 min natural release}
\subsection*{Ingredients}
\begin{itemize}
    \item Lean beef mince — 900 g
    \item Onion — 1 large, diced
    \item Garlic — 3 cloves, minced
    \item Canned tomatoes — 2 $\times$ 400 g
    \item Kidney beans — 2 $\times$ 400 g, drained and rinsed
    \item Black beans — 1 $\times$ 400 g, drained and rinsed
    \item Frozen mixed vegetables — 600 g
    \item Chilli powder — 2 tbsp
    \item Cumin — 1 tbsp
    \item Paprika — 1 tbsp
    \item Salt — 1$\tfrac{1}{4}$ tsp
    \item Stock or water — 1$\tfrac{1}{2}$ cups
    \item Oil — for sautéing
\end{itemize}

\subsection*{Method}
\begin{enumerate}
    \item Set the Instant Pot to \textbf{Sauté}. Add a small amount of oil and cook the onion for 2--3 minutes until softened.
    \item Add the beef mince and cook until well browned, breaking it up with a spoon.
    \item Stir in the garlic, chilli powder, cumin, and paprika. Cook for 30 seconds until fragrant.
    \item Add the canned tomatoes, kidney beans, black beans, frozen mixed vegetables, salt, and stock or water. Stir well, scraping the bottom to prevent sticking.
    \item Secure the lid and cook on \textbf{High Pressure} for 15 minutes.
    \item Allow a \textbf{Natural Pressure Release} for 10 minutes, then vent remaining pressure.
    \item Taste and adjust seasoning. If a thicker chilli is desired, use \textbf{Sauté} mode for a further 5 minutes.
\end{enumerate}

\subsection*{Serving Suggestions}
Serve with rice, baked potatoes, or tortillas. Also works well on its own for a lower-carb option. Suitable for meal prep and freezing.

\subsection*{Notes}
\begin{itemize}
    \item Frozen mixed vegetables add bulk, fibre, and nutrients at low cost with no additional prep.
    \item Using mixed beans stretches the mince while maintaining protein and satiety.
    \item This recipe freezes well for up to 3 months.
    \item Optional additions include frozen peas, frozen broccoli or cauliflower, or a handful of frozen spinach stirred through when reheating.
\end{itemize}
\end{recipebox}
